\textbf{Phần 1. Trắc nghiệm lựa chọn}
\setcounterpageref{ex}{0}

\begin{ex} %[4D8N45]
Ước lượng kết quả phép tính $49 \times 21$ bằng cách làm tròn các thừa số đến hàng chục ta được
\choice
{$800$}
{\True $1000$}
{$1200$}
{$900$}
\loigiai{
$50 \times 20 = 1000$.\\
\textbf{Đáp án của câu hỏi là B}
}
\end{ex}

\begin{ex} %[4D8N45]
Ước lượng kết quả của phép chia $398 : 19$ ta được khoảng
\choice
{\True $20$}
{$10$}
{$30$}
{$40$}
\loigiai{
$400 : 20 = 20$.\\
\textbf{Đáp án của câu hỏi là A}
}
\end{ex}

\begin{ex} %[4D8H45]
Ước lượng số tiền mua 19 quyển vở giá $9800$ đồng mỗi quyển bằng cách làm tròn ta được khoảng
\choice
{$180\,000$ đồng}
{$220\,000$ đồng}
{$150\,000$ đồng}
{\True $200\,000$ đồng}
\loigiai{
$20 \times 10\,000 = 200\,000$ đồng.\\
\textbf{Đáp án của câu hỏi là D}
}
\end{ex}

\begin{ex} %[4D8H45]
Một rạp chiếu phim có 28 hàng ghế, mỗi hàng có 32 ghế. Ước lượng tổng số ghế trong rạp là khoảng
\choice
{\True $900$ ghế}
{$600$ ghế}
{$1200$ ghế}
{$800$ ghế}
\loigiai{
$30 \times 30 = 900$ ghế.\\
\textbf{Đáp án của câu hỏi là A}
}
\end{ex}

\begin{ex} %[4D8H45]
Ước lượng chiều dài sân trường khoảng bao nhiêu mét?
\choice
{$5\text{ m}$}
{$500\text{ m}$}
{\True $50\text{ m}$}
{$5\text{ km}$}
\loigiai{
Chiều dài sân trường khoảng $40 - 60\text{ m}$.\\
\textbf{Đáp án của câu hỏi là C}
}
\end{ex}

\begin{ex} %[4D8H45]
Ước lượng kết quả tổng $4980 + 3020 + 2010$ khoảng
\choice
{$9000$}
{\True $10\,000$}
{$11\,000$}
{$8000$}
\loigiai{
$5000 + 3000 + 2000 = 10\,000$.\\
\textbf{Đáp án của câu hỏi là B}
}
\end{ex}

\begin{ex} %[4D8V45]
Mẹ mang $500\,000$ đồng đi siêu thị. Mẹ dự định mua: $1$ thùng sữa giá $295\,000$ đồng, $2\text{ kg}$ thịt lợn giá $110\,000$ đồng/kg. Mẹ có đủ tiền mua không?
\choice
{Vừa đủ}
{Thừa tiền rất nhiều}
{Không ước lượng được}
{\True Thiếu tiền (cần khoảng $515\,000$ đồng)}
\loigiai{
$295\,000 + 2 \times 110\,000 = 515\,000 > 500\,000$ đồng nên thiếu tiền.\\
\textbf{Đáp án của câu hỏi là D}
}
\end{ex}

\begin{ex} %[4D8V45]
Một người đi xe máy mỗi giờ đi được khoảng $38\text{ km}$. Ước lượng quãng đường người đó đi trong 5 giờ khoảng
\choice
{$150\text{ km}$}
{$250\text{ km}$}
{\True $200\text{ km}$}
{$300\text{ km}$}
\loigiai{
$40 \times 5 = 200\text{ km}$.\\
\textbf{Đáp án của câu hỏi là C}
}
\end{ex}

\begin{ex} %[4D8V45]
Một cửa hàng bán được $48$ chiếc áo thun với giá $195\,000$ đồng một chiếc. Ước lượng doanh thu bán áo khoảng
\choice
{$8\,000\,000$ đồng}
{$12\,000\,000$ đồng}
{$15\,000\,000$ đồng}
{\True $10\,000\,000$ đồng}
\loigiai{
$50 \times 200\,000 = 10\,000\,000$ đồng.\\
\textbf{Đáp án của câu hỏi là D}
}
\end{ex}

\begin{ex} %[4D8V45]
Một khối lớp 4 có 5 lớp, mỗi lớp có khoảng 32 học sinh. Ước lượng tổng số học sinh khối 4 là khoảng
\choice
{$120$ học sinh}
{$200$ học sinh}
{\True $150$ học sinh (hoặc 160 học sinh)}
{$180$ học sinh}
\loigiai{
$5 \times 30 = 150$ học sinh (chính xác là 160).\\
\textbf{Đáp án của câu hỏi là C}
}
\end{ex}

\begin{ex} %[4D8V45]
Một bể nước hình hộp chữ nhật chứa được khoảng $2950\text{ l}$ nước. Khi đã dùng hết một nửa, lượng nước còn lại khoảng
\choice
{\True $1500\text{ l}$}
{$1000\text{ l}$}
{$2000\text{ l}$}
{$2500\text{ l}$}
\loigiai{
$3000 : 2 = 1500\text{ l}$.\\
\textbf{Đáp án của câu hỏi là A}
}
\end{ex}

\begin{ex} %[4D8C45]
Để chuẩn bị nước uống cho chuyến dã ngoại 150 học sinh, biết mỗi học sinh uống khoảng $450\text{ ml}$ nước. Cần chuẩn bị khoảng bao nhiêu bình nước loại $20$ lít?
\choice
{$2$ bình}
{\True $4$ bình (khoảng $67{,}5\text{ l}$)}
{$6$ bình}
{$8$ bình}
\loigiai{
Tổng lượng nước: $150 \times 450 = 67\,500\text{ ml} = 67{,}5\text{ l}$.\\
Số bình 20 lít: $67{,}5 : 20 \approx 3{,}375 \Rightarrow$ cần 4 bình.\\
\textbf{Đáp án của câu hỏi là B}
}
\end{ex}

\textbf{Phần 2. Trắc nghiệm đúng sai}
\setcounterpageref{ex}{0}

\begin{ex} %[4D8H45]
Xét tính hợp lí của các ước lượng:
\choiceTF
{\True Cân nặng của một chiếc xe máy khoảng $100\text{ kg}$.}
{Thời gian ngủ mỗi đêm của học sinh tiểu học khoảng $15$ giờ.}
{\True Thời gian ngủ mỗi đêm của học sinh tiểu học khoảng $8 - 9$ giờ.}
{\True Ước lượng tích $98 \times 49$ khoảng $100 \times 50 = 5000$.}
\loigiai{
\begin{itemize}
    \item Xe máy ~100 kg: Đúng.
    \item Học sinh ngủ 8-9h (15h là sai). Ý 2 Sai, Ý 3 Đúng.
    \item Ước lượng tích đúng.
\end{itemize}
\begin{itemize}
    \item a) \textbf{Đúng}.
    \item b) \textbf{Sai}.
    \item c) \textbf{Đúng}.
    \item d) \textbf{Đúng}.
\end{itemize}
}
\end{ex}

\begin{ex} %[4D8V45]
Bạn An có $100\,000$ đồng đi mua đồ dùng học tập: mua 1 hộp bút giá $38\,000$ đồng, 5 quyển vở giá $9500$ đồng/quyển và 1 cây thước kẻ giá $5000$ đồng. Xét tính đúng sai:
\choiceTF
{\True Tiền mua 5 quyển vở hết $47\,500$ đồng.}
{\True Tổng số tiền mua 3 món đồ là $90\,500$ đồng.}
{Bạn An không đủ tiền để thanh toán.}
{\True Bạn An còn thừa khoảng $9500$ đồng.}
\loigiai{
\begin{itemize}
    \item Vở: $5 \times 9500 = 47\,500$ đồng: Đúng.
    \item Tổng: $38\,000 + 47\,500 + 5000 = 90\,500$ đồng: Đúng.
    \item Đủ tiền vì $90\,500 < 100\,000$. Sai.
    \item Thừa: $100\,000 - 90\,500 = 9500$ đồng. Đúng.
\end{itemize}
\begin{itemize}
    \item a) \textbf{Đúng}.
    \item b) \textbf{Đúng}.
    \item c) \textbf{Sai}.
    \item d) \textbf{Đúng}.
\end{itemize}
}
\end{ex}

\textbf{Phần 3. Trả lời ngắn}
\setcounterpageref{ex}{0}

\begin{ex} %[4D8H45]
Ước lượng tích $51 \times 39$ bằng cách làm tròn đến hàng chục.
\shortans{2000}
\loigiai{
\textbf{Đáp số:} $2000$\\
$50 \times 40 = 2000$.
}
\end{ex}

\begin{ex} %[4D8H45]
Ước lượng thương $595 : 31$ bằng cách làm tròn.
\shortans{20}
\loigiai{
\textbf{Đáp số:} $20$\\
$600 : 30 = 20$.
}
\end{ex}

\begin{ex} %[4D8V45]
Một cửa hàng bán được 29 cái cặp sách giá $190\,000$ đồng một cái. Ước lượng doanh thu theo triệu đồng (làm tròn đến hàng triệu).
\shortans{6}
\loigiai{
\textbf{Đáp số:} $6$\\
$30 \times 200\,000 = 6\,000\,000$ đồng $= 6$ triệu đồng.
}
\end{ex}

\begin{ex} %[4D8V45]
Một xe tải chở 48 thùng hàng, mỗi thùng nặng $21\text{ kg}$. Ước lượng tổng khối lượng hàng theo ki-lô-gam (làm tròn đến hàng trăm).
\shortans{1000}
\loigiai{
\textbf{Đáp số:} $1000$\\
$50 \times 20 = 1000\text{ kg}$.
}
\end{ex}

\begin{ex} %[4D8V45]
Ước lượng tổng $1980 + 2030 + 4990$ (làm tròn đến hàng nghìn).
\shortans{9000}
\loigiai{
\textbf{Đáp số:} $9000$\\
$2000 + 2000 + 5000 = 9000$.
}
\end{ex}

\begin{ex} %[4D8C45]
Một trường học có 980 học sinh, mỗi ngày mỗi học sinh uống khoảng $250\text{ ml}$ sữa tươi. Ước lượng số lít sữa toàn trường tiêu thụ mỗi ngày (làm tròn đến hàng chục lít).
\shortans{250}
\loigiai{
\textbf{Đáp số:} $250$\\
$1000 \times 250\text{ ml} = 250\,000\text{ ml} = 250$ lít.
}
\end{ex}
